
% ---------------------------------------------------------------------------
% Case 1: COMPAS
% ---------------------------------------------------------------------------
\begin{frame}{Case 1: COMPAS Recidivism Scores (2016)}

    \textbf{The system.} COMPAS is a commercial risk-assessment tool (Northpointe, now Equivant)
    that scores a criminal defendant's likelihood of re-offending on a 1--10 scale. US courts have
    used such scores to inform bail, sentencing, and parole decisions.

    \vspace{0.8em}

    \textbf{The analysis.} ProPublica journalists obtained scores for more than 7{,}000 people
    arrested in Broward County, Florida in 2013--2014 and checked, two years later, who actually
    re-offended.

    \vspace{0.8em}

    \begin{columns}[T]
        \begin{column}{0.52\textwidth}
            \centering
            \small
            \begin{tabular}{p{3.5cm}cc}
                \toprule
                & \textbf{White} & \textbf{Black} \\
                \midrule
                Labelled higher risk, but did \emph{not} re-offend & 23.5\% & 44.9\% \\
                \addlinespace
                Labelled lower risk, yet \emph{did} re-offend & 47.7\% & 28.0\% \\
                \bottomrule
            \end{tabular}
        \end{column}
        \begin{column}{0.44\textwidth}
            \begin{itemize}\small
                \item Overall, 61\% of those scored likely to re-offend were arrested for some
                      crime within two years.
                \item Of those predicted to commit a \emph{violent} crime, 20\% did.
            \end{itemize}
        \end{column}
    \end{columns}

    \vspace{0.5em}
    {\scriptsize J.~Angwin, J.~Larson, S.~Mattu and L.~Kirchner, ``Machine Bias,'' \emph{ProPublica},
    23 May 2016, \href{https://www.propublica.org/article/machine-bias-risk-assessments-in-criminal-sentencing}{propublica.org}.}
\end{frame}
